\section{Results}\label{sec:results}

All three scenarios have completed discrimination reports for repetitions 1--5: 11 models for each KFRE-input scenario and 10 for the twenty-feature scenario, yielding 160 model--scenario--run records. Tables~\ref{tab:discrimination} and~\ref{tab:twenty} summarize their reported metrics as mean (sample SD). These runs reuse a fixed patient partition; their variation reflects repeated model fitting and tuning, not independent test cohorts or repeated random splits. Rep99 feasibility results are excluded. The evaluation qualifications in Section~\ref{sec:eval} apply throughout.

\subsection{Four and Eight Features}

\begin{table}[t]
\centering
\footnotesize
\setlength{\tabcolsep}{3pt}
\begin{tabular}{@{}lccc@{}}
\toprule
\multicolumn{4}{c}{\textbf{Four features}} \\
\textbf{Model} & \textbf{C-index}$\uparrow$ & \textbf{IBS}$\downarrow$ & \textbf{Mean AUC}$\uparrow$ \\
\midrule
Cox & 0.308 (0.000) & 0.531 (0.000) & 0.272 (0.000) \\
Dynamic-DeepHit & \textbf{0.827 (0.011)} & \textbf{0.279 (0.018)} & \textbf{0.939 (0.017)} \\
Hazard Transformer & 0.637 (0.001) & 0.409 (0.007) & 0.701 (0.002) \\
Logistic Hazard & 0.551 (0.012) & 0.416 (0.006) & 0.567 (0.018) \\
RNN-Surv & 0.522 (0.006) & 0.430 (0.011) & 0.526 (0.011) \\
DeepSurv & 0.540 (0.020) & 0.463 (0.075) & 0.549 (0.023) \\
GBSA & 0.518 (0.003) & 0.438 (0.008) & 0.530 (0.008) \\
Survival RF & 0.532 (0.005) & 0.445 (0.003) & 0.543 (0.007) \\
Survival SVM & 0.543 (0.000) & 0.417 (0.000) & 0.571 (0.000) \\
Weibull AFT & 0.552 (0.000) & 0.414 (0.000) & 0.579 (0.000) \\
KFRE & 0.539 (0.000) & 0.552 (0.000) & 0.553 (0.000) \\
\midrule
\multicolumn{4}{c}{\textbf{Eight features}} \\
\textbf{Model} & \textbf{C-index}$\uparrow$ & \textbf{IBS}$\downarrow$ & \textbf{Mean AUC}$\uparrow$ \\
\midrule
Cox & 0.470 (0.000) & 0.578 (0.000) & 0.425 (0.000) \\
Dynamic-DeepHit & \textbf{0.880 (0.009)} & \textbf{0.281 (0.024)} & \textbf{0.953 (0.009)} \\
Hazard Transformer & 0.605 (0.039) & 0.529 (0.205) & 0.631 (0.049) \\
Logistic Hazard & 0.521 (0.019) & 0.490 (0.016) & 0.500 (0.029) \\
RNN-Surv & 0.531 (0.005) & 0.493 (0.005) & 0.590 (0.015) \\
DeepSurv & 0.484 (0.014) & 0.555 (0.115) & 0.486 (0.022) \\
GBSA & 0.541 (0.011) & 0.529 (0.014) & 0.559 (0.011) \\
Survival RF & 0.522 (0.008) & 0.523 (0.007) & 0.539 (0.008) \\
Survival SVM & 0.540 (0.000) & 0.488 (0.000) & 0.595 (0.000) \\
Weibull AFT & 0.554 (0.000) & 0.485 (0.000) & 0.597 (0.000) \\
KFRE & 0.539 (0.000) & 0.737 (0.000) & 0.579 (0.000) \\
\bottomrule
\end{tabular}
\caption{Mean (sample SD) over five runs on a fixed patient partition. C-index uses available follow-up; IBS integrates days 1--730 after the common risk-score transformation; AUC is averaged over the daily grid before day 730. Bold marks the best recorded mean in each scenario. SD describes run variation, not sampling uncertainty; 0.000 denotes rounding at the displayed precision.}
\label{tab:discrimination}
\end{table}

\textbf{Dynamic-DeepHit has the highest recorded discrimination in both scenarios.} Its mean C-index is 0.827 (SD 0.011) with four features and 0.880 (0.009) with eight features; the corresponding mean time-dependent AUCs are 0.939 (0.017) and 0.953 (0.009). KFRE has C-index 0.539 in both scenarios, with AUC 0.553 and 0.579, respectively. Dynamic-DeepHit also has the lowest transformed-score IBS, 0.279 (0.018) and 0.281 (0.024). This IBS is computed from the shared scalar-risk-to-survival transformation and does not establish that its native probability predictions are the best calibrated.

Hazard Transformer has the next highest mean C-index, 0.637 (0.001) and 0.605 (0.039). The remaining learned models, excluding Cox, span mean C-indices of 0.518--0.552 with four features and 0.484--0.554 with eight features. These are descriptive comparisons: the run SDs do not provide confidence intervals for differences from KFRE or establish equivalence. Cox remains below 0.5, at 0.308 and 0.470, and DeepSurv also has a mean below 0.5 in the eight-feature scenario.

\textbf{Feature-set differences are model-dependent.} Moving from four to eight features changes Dynamic-DeepHit's mean C-index by +0.053, Cox's by +0.162, DeepSurv's by -0.056, and Hazard Transformer's by -0.032. The other seven models change by 0.03 or less. The complete-case cohorts differ in size and composition (2{,}809 versus 1{,}213 patients), so these changes cannot be attributed to the four additional measurements alone.

\subsection{Full Twenty-Feature Cohort}

\begin{table}[t]
\centering
\footnotesize
\setlength{\tabcolsep}{3pt}
\begin{tabular}{@{}lccc@{}}
\toprule
\multicolumn{4}{c}{\textbf{Twenty features (heterogeneous)}} \\
\textbf{Model} & \textbf{C-index}$\uparrow$ & \textbf{IBS}$\downarrow$ & \textbf{Mean AUC}$\uparrow$ \\
\midrule
Cox & 0.457 (0.000) & 0.397 (0.000) & 0.433 (0.000) \\
Dynamic-DeepHit & \textbf{0.856 (0.035)} & \textbf{0.248 (0.013)} & \textbf{0.946 (0.026)} \\
Hazard Transformer & 0.717 (0.060) & 0.314 (0.032) & 0.802 (0.078) \\
Logistic Hazard & 0.522 (0.005) & 0.374 (0.006) & 0.531 (0.026) \\
RNN-Surv & 0.523 (0.001) & 0.376 (0.003) & 0.527 (0.001) \\
DeepSurv & 0.533 (0.003) & 0.370 (0.001) & 0.539 (0.007) \\
GBSA & 0.521 (0.000) & 0.371 (0.000) & 0.529 (0.000) \\
Survival RF & 0.522 (0.001) & 0.372 (0.000) & 0.534 (0.001) \\
Survival SVM & 0.513 (0.000) & 0.369 (0.000) & 0.540 (0.000) \\
Weibull AFT & 0.525 (0.000) & 0.370 (0.000) & 0.546 (0.000) \\
\bottomrule
\end{tabular}
\caption{Full-cohort twenty-feature results, mean (sample SD), five runs. The cohort contains 32{,}601 patients, with 6{,}521 held out. Metric definitions and variability interpretation follow Table~\ref{tab:discrimination}; KFRE has no corresponding equation.}
\label{tab:twenty}
\end{table}

The heterogeneous scenario now reports the full extracted cohort of 32{,}601 patients and approximately 8.13 million lab-event rows, replacing the earlier 20-patient feasibility analysis. Dynamic-DeepHit has the highest mean C-index, 0.856 (0.035), and mean AUC, 0.946 (0.026), with IBS 0.248 (0.013). Hazard Transformer follows at C-index 0.717 (0.060) and AUC 0.802 (0.078); its poor pilot result does not describe the full-cohort runs. Both sequence-based architectures vary more across runs here than Dynamic-DeepHit does in either complete-case scenario. The other non-Cox models have mean C-indices between 0.513 and 0.533, while Cox remains below chance at 0.457.

These results characterize a different cohort and input representation, with 20 laboratory values and their missingness flags. They are not a matched-patient experiment establishing the isolated effect of increasing feature count.

\subsection{Feature Importance}

Figure~\ref{fig:feature_importance_four} illustrates the updated four-feature analysis for rep1. Importance measures differ by model: absolute coefficients for Cox, Weibull AFT, and Survival SVM; native feature importance for GBSA; permutation importance for Survival RF; and gradient-based summaries for neural models. Magnitudes are not comparable across these families.

\begin{figure}[t]
\centering
\includegraphics[width=\textwidth]{figs/four_features_all_models_feature_importance.png}
\caption{Four-feature importance diagnostics from rep1, regenerated with the completed analyses. Panels use model-specific measures and are illustrative of one run; their scales are not comparable. These are not SHAP values.}
\label{fig:feature_importance_four}
\end{figure}

In this illustrative run, eGFR has the largest reported neural importance for Dynamic-DeepHit (0.293), Hazard Transformer (0.463), Logistic Hazard (0.402), and RNN-Surv (0.382). GBSA assigns similar importance to uACR (0.386) and eGFR (0.376), while DeepSurv's largest value is for sex (0.277). These within-model summaries do not establish a common biological ranking. In particular, a raw uACR coefficient displayed as 0.0000 is scale- and precision-dependent and does not demonstrate that uACR was discarded.

All five twenty-feature importance reports record memory-allocation failures for the Dynamic-DeepHit and Hazard Transformer gradient calculations. The resulting uniform Dynamic-DeepHit values and first-layer-weight Transformer fallbacks are excluded from substantive feature-importance conclusions.

\subsection{Calibration and Decision-Curve Diagnostics}

Calibration is assessed separately from the transformed-score IBS, using native horizon-specific probabilities where available and a training-fitted approximation otherwise (Section~\ref{sec:eval}). Figure~\ref{fig:calibration_four} gives the updated four-feature rep1 diagnostic. These plots are descriptive and do not provide calibration confidence intervals.

\begin{figure}[t]
\centering
\includegraphics[width=\textwidth]{figs/four_features_calibration_plot.png}
\caption{Four-feature calibration diagnostics from rep1 at 2 and 5 years: mean predicted risk versus Kaplan--Meier-observed risk within risk groups. Native and approximated probabilities are distinguished in the analysis reports. This figure illustrates one run, not an average over five runs.}
\label{fig:calibration_four}
\end{figure}

\textbf{High discrimination does not guarantee stable absolute-risk calibration.} Dynamic-DeepHit's four-feature rep1 predictions broadly track the observed risks, but calibration varies across scenarios and runs. In the twenty-feature rep1 analysis, its 2-year risk-group means lie between 0.166 and 0.187 while the corresponding observed risks range from 0.431 to 0.998; at 5 years, predictions range from 0.390 to 0.423 versus observed risks of 0.676--1.000. This substantial underprediction precludes a general claim of uniformly good calibration. Hazard Transformer also produces constant or near-constant native probabilities at the horizon in multiple reports, even where its scalar risk score discriminates between patients.

\textbf{The current decision curves do not support a clinical-utility comparison.} The saved model curves use patient-level predictions, whereas the treat-all and eGFR-rule comparators use raw test rows (and, for the heterogeneous eGFR rule, a measured-eGFR subset). The eGFR rule also derives its displayed threshold from the fraction referred rather than using each common decision threshold. These differing evaluation units and thresholds invalidate direct net-benefit comparisons. We therefore omit the decision-curve figures and make no claim that model-guided referral improves on treat-all or an eGFR cutoff. Harmonized patient-level recomputation is required.
